\section{Eksperymenty}
\label{sec:eksperymenty}

W~bieżącej wersji projektu wdrożony jest mechanizm zbierania metryk
populacji (\texttt{print\_population\_stats} w~\texttt{initial\_population.py})
oraz walidator \texttt{examples/validate\_population.py}, co stanowi
podstawę pod planowane eksperymenty. Pełna seria pomiarów porównawczych
nie została jeszcze przeprowadzona --- w~tym rozdziale prezentujemy
jedynie zarys planu eksperymentalnego oraz dane, jakie są możliwe do
uzyskania z~aktualnego kodu.

\subsection{Plan eksperymentalny}

\paragraph{Zmienne niezależne.}
\begin{itemize}
    \item Rozmiar mapy: $H \times W \in \{10\!\times\!10,\ 20\!\times\!20,\ 30\!\times\!30\}$.
    \item Gęstość atrakcji: liczba atrakcji $\in \{4, 15, 30\}$.
    \item Rozkład wag: \texttt{uniform} vs.\ \texttt{clusters}.
    \item Parametry ABC: $N \in \{20, 50, 100\}$, $I \in \{50, 200, 500\}$,
          $L \in \{30, 60, 100\}$, $O \in \{N/2, N, 2N\}$.
    \item Ziarno losowości \texttt{seed} --- minimum 5 powtórzeń dla każdej
          kombinacji.
\end{itemize}

\paragraph{Zmienne zależne (mierzone).}
\begin{itemize}
    \item Frakcja rozwiązań dopuszczalnych w~populacji końcowej.
    \item Średnia, mediana i~maksimum \texttt{total\_value} w~populacji.
    \item Średni \texttt{movement\_cost}, \texttt{total\_time}, \texttt{attraction\_cost}.
    \item Średnia długość ścieżki (liczba pól).
    \item Liczba unikalnych ścieżek (różnorodność populacji).
    \item Czas wykonania (sekundy) jako funkcja $N$, $I$, $O$.
\end{itemize}

\paragraph{Hipotezy do weryfikacji.}
\begin{enumerate}
    \item Zwiększenie $I$ (liczby iteracji) podnosi medianę \texttt{total\_value},
          ale z~malejącym przyrostem (krzywa nasycenia).
    \item Faza onlookerów wzmacnia eksploatację --- wyższe $O$ przy stałym $I$
          poprawia średnią wartość celu, ale zmniejsza różnorodność populacji.
    \item Wyższy limit stagnacji $L$ zwiększa eksploatację kosztem
          eksploracji; bardzo niskie $L$ powoduje nadmierne ,,resetowanie''
          przez zwiadowców i~uniemożliwia dopracowanie rozwiązań.
    \item Rozkład wag \texttt{clusters} (klastry niskokosztowe) daje wyższe
          wartości \texttt{total\_value} w~tym samym budżecie czasu niż
          \texttt{uniform}.
\end{enumerate}

\subsection{Porównanie z baselineami}

W~kodzie dostępne są dwa proste algorytmy referencyjne:
\texttt{solve\_random\_walk} i~\texttt{solve\_greedy\_attractions}
(\texttt{src/path\_solver.py}). Porównawcza tabela ma postać:

\begin{table}[H]
\centering
\caption{Szablon zestawienia wyników (do uzupełnienia).}
\label{tab:results-template}
\begin{tabular}{lcccc}
\toprule
Algorytm & \texttt{value} & \texttt{move\_cost} & \texttt{total\_time} & \texttt{feasible} \\
\midrule
Random Walk        & -- & -- & -- & -- \\
Greedy attractions & -- & -- & -- & -- \\
ABC (best in pop)  & -- & -- & -- & -- \\
ABC (mean in pop)  & -- & -- & -- & -- \\
\bottomrule
\end{tabular}
\end{table}

\subsection{Zbieranie wyników}

Aktualny stan kodu pozwala uzyskać następujące dane bez modyfikacji
implementacji:

\begin{itemize}
    \item Surowe statystyki (\Cref{lst:stats} z~\Cref{sec:aplikacja}) ---
          wypisywane na stdout po uruchomieniu \texttt{initial\_population.py}.
    \item Plik \texttt{initial\_population.json} z~wszystkimi rozwiązaniami
          --- pozwala policzyć dowolne dodatkowe statystyki ex post
          (np.\ w~notatniku).
    \item Wizualizację (\texttt{visualize\_paths.py}) dla wybranych \texttt{id}.
\end{itemize}

\subsection{Wyniki wstępne}

Pełne wyniki eksperymentalne zostaną dołączone w~kolejnej wersji
dokumentacji. Bieżąca implementacja przeszła testy poprawności:
\begin{itemize}
    \item Wszystkie wygenerowane rozwiązania w~populacji końcowej spełniają
          wszystkie ograniczenia (\texttt{feasible=true}) --- co jest
          gwarantowane konstrukcyjnie przez algorytm.
    \item Walidator \texttt{validate\_population.py} potwierdza brak
          powtórzeń krawędzi w~każdej trasie.
    \item Wizualizacja prezentuje sensowne kształty tras (krzywe biegnące
          przez atrakcje, omijające zablokowane atrakcje niezaplanowane).
\end{itemize}
